\documentclass[12pt,a4paper]{article}
\usepackage[utf8]{inputenc}
\usepackage[T1]{fontenc}
\usepackage[brazil]{babel}
\usepackage{amsmath,amssymb}
\usepackage{geometry}
\geometry{margin=2.5cm}
\title{Material de Estudo: Resultado Final --- Fechamento da Disciplina}
\author{Princ\'ipio de Modelagem Matem\'atica}
\date{Unidade 36}
\begin{document}
\maketitle

\section*{Objetivo}
Revisar os componentes da nota final, consolidar as compet\^encias desenvolvidas e apontar perspectivas de pesquisa.

\section*{Componentes da Nota Final}

\begin{tabular}{|l|c|}
\hline
\textbf{Componente} & \textbf{Peso} \\
\hline
Prova I (Unidades 1--15) & 25\% \\
Prova II (Unidades 16--30) & 25\% \\
Trabalho Final (relat\'orio + c\'odigo) & 30\% \\
Apresenta\c{c}\~ao Oral & 15\% \\
Participa\c{c}\~ao (exerc\'icios, avalia\c{c}\~ao por pares) & 5\% \\
\hline
\textbf{Total} & 100\% \\
\hline
\end{tabular}

\section*{O Que um Modelador Deve Saber}

Ao concluir esta disciplina, voc\^e deve ser capaz de:
\begin{enumerate}
  \item \textbf{Formular} um modelo matem\'atico a partir de um problema real.
  \item \textbf{Classificar} o modelo (determin\'istico/estoc\'astico, discreto/cont\'inuo, anal\'itico/num\'erico).
  \item \textbf{Resolver} ou \textbf{simular} o modelo com m\'etodos apropriados.
  \item \textbf{Validar} os resultados contra dados ou solu\c{c}\~oes conhecidas.
  \item \textbf{Interpretar} os resultados em termos do fen\^omeno original.
  \item \textbf{Comunicar} os resultados com clareza e rigor.
  \item \textbf{Reconhecer} as limita\c{c}\~oes do modelo e propor melhorias.
\end{enumerate}

\section*{Ciclo Completo da Modelagem}

\[
\text{Problema} \xrightarrow{\text{formular}} \text{Modelo} \xrightarrow{\text{resolver/simular}} \text{Resultados} \xrightarrow{\text{validar}} \text{Interpreta\c{c}\~ao} \xrightarrow{\text{refinar}} \text{Modelo}
\]

Este ciclo \'e iterativo e nunca termina --- todo modelo \'e uma aproxima\c{c}\~ao.

\section*{Perspectivas de Pesquisa}

Os m\'etodos da disciplina s\~ao ativamente usados em:
\begin{itemize}
  \item \textbf{Biologia e sa\'ude:} modelos de epidemias (COVID-19), din\^amica de prote\'inas, redes neuronais.
  \item \textbf{Engenharia de transportes:} simula\c{c}\~ao de tr\'afego urbano, redes de mobilidade.
  \item \textbf{F\'isica e qu\'imica computacional:} DM de materiais, l\'iquidos, biomoléculas.
  \item \textbf{Ci\^encias sociais:} modelos baseados em agentes, forma\c{c}\~ao de opini\~ao.
  \item \textbf{Ecologia:} din\^amica de popula\c{c}\~oes, modelos de extinc\~ao.
  \item \textbf{Ci\^encia de dados:} regressão, aprendizado de m\'aquina (continua\c{c}\~ao natural).
\end{itemize}

\section*{Exerc\'icios Finais}

\begin{enumerate}
  \item \textbf{(Autoavalia\c{c}\~ao)} Para cada compet\^encia da lista acima, d\^e uma nota de 1--5 para si mesmo. Quais s\~ao seus pontos fortes? Quais precisam de desenvolvimento?
  
  \item \textbf{(Mapa mental)} Crie um mapa mental conectando os 36 t\'opicos da disciplina. Quais os n\'os centrais (conceitos que aparecem em v\'arios t\'opicos)?
  
  \item \textbf{(Modelo favorito)} Escreva um par\'agrafo sobre o modelo que mais lhe interessou. O que o torna elegante? Quais s\~ao suas applica\c{c}\~oes?
  
  \item \textbf{(Frase de Box)} George Box disse: ``Todos os modelos s\~ao errados, mas alguns s\~ao \'uteis.'' Interprete essa afirma\c{c}\~ao com um exemplo concreto da disciplina. O que torna um modelo ``\'util'' mesmo sendo ``errado''?
  
  \item \textbf{(Projeto futuro)} Elabore uma proposta de projeto de pesquisa de 1 p\'agina usando at least dois m\'etodos diferentes da disciplina. Inclua: motiva\c{c}\~ao, pergunta, m\'etodos, resultados esperados, desafios.
  
  \item \textbf{(Conex\~ao)} Conecte a disciplina com sua \'area de estudo principal. Que problema relevante em sua \'area poderia ser abordado com modelagem matem\'atica?
  
  \item \textbf{(Legado)} O que voc\^e levaria desta disciplina para sua vida profissional e acad\^emica? Quais h\'abitos de pensamento ou ferramentas s\~ao mais valiosos?
\end{enumerate}

\section*{Mensagem de Encerramento}

\begin{quote}
\textit{``Todos os modelos s\~ao errados, mas alguns s\~ao \'uteis.''}\\
--- George E.\ P.\ Box (1919--2013)
\end{quote}

A modelagem matem\'atica n\~ao \'e sobre encontrar a verdade absoluta, mas sobre construir representa\c{c}\~oes \'uteis da realidade que nos permitem compreender, prever e controlar fen\^omenos. A arte est\'a em escolher as simplifica\c{c}\~oes certas --- o bastante para capturar o essencial, mas n\~ao tanto que se perca o fen\^omeno.

Voc\^e agora possui um arsenal: an\'alise dimensional, expans\~oes de Taylor, EDOs, EDEs, cadeias de Markov, RNG, aut\^omatos celulares, m\'etodos num\'ericos e din\^amica molecular. Use-os com rigor, criatividade e honestidade cient\'ifica.

\medskip
\noindent\textbf{Bom trabalho a todos!}

\section*{Resumo Final da Disciplina}

\begin{tabular}{|l|l|}
\hline
\textbf{Bloco} & \textbf{T\'opicos} \\
\hline
Fundamentos (1--8) & An\'alise dim., escalas, fractais, Taylor, erros \\
Modelos anal\'iticos (9--15) & Regressão, EDO, exponencial, tr\'afego, Prova I \\
Estocasticidade (16--20) & EDE, Markov, RNG, caminhante aleat\'orio \\
M\'etodos num\'ericos (21--27) & CA, differ.\ finitas, RK4, DM I \\
DM e aplica\c{c}\~oes (28--31) & DM II-III, discuss\~ao, Prova II \\
Apresenta\c{c}\~oes (32--36) & Comunica\c{c}\~ao, avalia\c{c}\~ao, s\'intese, resultado \\
\hline
\end{tabular}

\section*{Refer\^encias Gerais}
\begin{itemize}
  \item Strogatz, S.\ H.\ \textit{Nonlinear Dynamics and Chaos}. Westview, 2015.
  \item Frenkel, D.; Smit, B.\ \textit{Understanding Molecular Simulation}. Academic Press, 2002.
  \item Wolfram, S.\ \textit{A New Kind of Science}. Wolfram Media, 2002.
  \item Press, W.\ H.\ et al.\ \textit{Numerical Recipes}. Cambridge, 2007.
  \item Notas de aula da disciplina.
\end{itemize}

\end{document}
